% Bewertungspaket zum Gesamttest «Planimetrie» — Quelle des PDFs (python3 scripts/build-lp-pdf.py planimetrie).
% Ganze Punkte je Teilschritt; Werte mit python3 nachgerechnet (06.10.2026, Folgefehler-Fälle eingeschlossen).
\documentclass[11pt]{article}
\usepackage{../lp-druck}
\renewcommand{\lpname}{Planimetrie}
\renewcommand{\lpbereich}{Grundlagenfach}
\renewcommand{\lpteilgebiet}{5.2}
\renewcommand{\lpfuss}{Bewertungspaket}
\newcolumntype{P}{>{\centering\arraybackslash}p{10mm}}
\newenvironment{raster}{\par\smallskip\noindent\tabularx{\linewidth}{@{}>{\raggedright\arraybackslash}p{38mm}P X@{}}\toprule
  \small\textsc{Teilschritt} & \small\textsc{P} & \small\textsc{Musterlösung}\\\midrule}{\bottomrule\endtabularx\par}
\newcommand{\fehler}[1]{\par\smallskip{\small\textcolor{lprot}{\bfseries Typische Fehler:} #1}}
\newcommand{\E}{\,{\footnotesize\bfseries\color{lpgruen}(E)}}
\newcommand{\A}{\,{\footnotesize\bfseries\color{lpblau}(A)}}
\begin{document}
\lpkopf{Bewertungspaket zum Gesamttest}{}

\begin{lpachtung}{Erst lösen, dann öffnen}
Dieses Paket enthält die Musterlösung. Wer es vor dem Lösen liest, misst am Ende nur, wie gut das Abschreiben gelingt.
\end{lpachtung}

\section*{1 · So gehst du vor}
\begin{enumerate}[leftmargin=*,itemsep=2pt]
\item \textbf{Gesamttest lösen} — vollständig, mit Skizze und Rechenweg, ohne dieses Paket; der Taschenrechner ist erlaubt.
\item \textbf{Fotografieren oder scannen}, gut lesbar, eine Seite pro Bild. \textbf{Kein Name} auf den Bildern.
  Handyfotos tragen oft unsichtbar Standort, Zeit und Gerät mit (Metadaten). Lade darum lieber einen Scan oder ein
  Bildschirmfoto des Fotos hoch, oder entferne vorher die Standortangaben.
\item \textbf{Bewerten lassen.} Ob du dafür eine KI nutzt und welche, entscheidest du selbst und in eigener Verantwortung —
  je nachdem, was dir aufgrund deines Alters und deiner persönlichen Situation erlaubt ist (Altersgrenzen und
  Nutzungsbedingungen des Anbieters, allenfalls Einverständnis deiner Eltern). Lade nur deine Lösung und dieses PDF hoch,
  keine persönlichen Daten, und füge den Auftrag aus Abschnitt~2 ein. Ohne KI kannst du dich mit Abschnitt~3 selbst bewerten.
\item \textbf{Rückmeldung prüfen.} Stimmt, was die KI aus deiner Handschrift gelesen hat? Eine KI kann sich irren —
  im Zweifel gilt das Raster in Abschnitt~3.
\item \textbf{Gezielt wiederholen} nach Abschnitt~4.
\end{enumerate}

\needspace{14\baselineskip}
\section*{2 · Auftrag an die KI}
\begin{tcolorbox}[colback=lplinie!25,colframe=lplinie,boxrule=0.5pt,arc=1mm,fontupper=\small]
Du bist Korrektorin oder Korrektor für Mathematik an einer Schweizer Berufsmaturitätsschule. Im Anhang findest du (1)~das Bewertungspaket zum Gesamttest «Planimetrie» mit Musterlösung und Punkteraster und (2)~Fotos meiner handschriftlichen Lösung.\par\medskip
Geh so vor:\par
1. Schreib zu jeder Aufgabe G1 bis G7 zuerst ab, was du in meiner Lösung liest — Rechenweg und Ergebnis. Was du nicht sicher lesen kannst, markierst du mit [unsicher]. Nicht raten.\par
2. Vergib die Punkte genau nach dem Raster im Bewertungspaket (Abschnitt 3), ganze Punkte je Zeile. Ziehe einen Fehler nur einmal ab: Rechne ich mit einem falschen Zwischenergebnis richtig weiter, gibt es die Folgepunkte. Ausnahme: Zeilen mit \textbf{(E)} gibt es nur für das richtige Ergebnis, nie als Folgepunkt. Die Zeilen «Typische Fehler» sagen, welche Rasterzeile bei diesem Fehler 0 Punkte gibt — sie sind kein zusätzlicher Abzug.\par
3. Andere richtige Lösungswege zählen voll, ebenso gleichwertige Schreibweisen: Dezimalpunkt oder Dezimalkomma, exakte Werte wie $6\pi$ oder $\sqrt{45}$ statt Dezimalzahlen, anders, aber vernünftig gerundete Ergebnisse (Abweichung höchstens $\pm 0.01$ in der letzten Stelle oder durch Runden von Zwischenergebnissen erklärbar), $\text{cm}^2$ statt $\text{m}^2$, wenn richtig umgerechnet. \textbf{Einheiten:} Fehlt bei einem Endergebnis die Einheit oder ist sie falsch ($\text{cm}$ statt $\text{cm}^2$), gibt es für das ganze Gesamttest einmal 1 Punkt Abzug, nicht je Aufgabe. Drei Stufen: Zeilen mit \textbf{(A)} verlangen nur Ablesen, diese Punkte gibt es auch ohne Rechenweg. Zeilen mit \textbf{(E)} sind Ergebnispunkte: Das Ergebnis muss stimmen \emph{und} aus einem erkennbaren Weg hervorgehen, ausser die Zeile sagt ausdrücklich etwas anderes. Alle übrigen Zeilen bewerten den Weg selbst.\par
4. Begründe jeden Abzug in einem Satz und nenne die Fehlerart (zum Beispiel «Seite statt Höhe» oder «Durchmesser statt Radius»). Schreib die Musterlösung nicht ab — zeig mir, wo mein Fehler liegt.\par
5. Gib zum Schluss eine Tabelle «Aufgabe | Punkte | Begründung», die Gesamtpunktzahl (von 25) und — nach der Selbsteinschätzung in Abschnitt 4 — welches Kapitel des Leitprogramms ich wiederholen soll. Keine Note.\par
6. Fehlt ein Foto oder ist eine Aufgabe unlesbar, sag es, statt sie zu bewerten.
\end{tcolorbox}

\needspace{16\baselineskip}
\section*{3 · Musterlösung und Punkteraster}
\textbf{Allgemein:} Ganze Punkte je Zeile. Ein Fehler wird nur einmal abgezogen (Folgefehler: wer mit einem falschen
Zwischenergebnis richtig weiterrechnet, bekommt die Wegpunkte der folgenden Zeilen, aber keine (E)-Punkte). In der Spalte P:
\textbf{(A)}~= nur ablesen, zählt auch ohne Rechenweg · \textbf{(E)}~= Ergebnispunkt, Ergebnis richtig
\emph{und} Weg erkennbar (ausser die Zeile sagt ausdrücklich etwas anderes) · ohne Zeichen = die Zeile bewertet den Weg.
Taschenrechner erlaubt; Ergebnisse auf zwei Dezimalen, Toleranz $\pm 0.01$. \textbf{Einheiten:} fehlende oder falsche
Einheit einmal 1 Punkt Abzug für den ganzen Test. «Typische Fehler» nennt, welche Zeile 0 Punkte gibt, und die Punkte, die bleiben.
\textbf{Teilrichtige Zeilen:} Verlangt eine Zeile mehrere Angaben und stimmt nur ein Teil, gibt die Zeile 0 Punkte.

\begin{lpaufgabe}{G1}{Dreiecke beschreiben}{3 P}
\begin{raster}
(a) Begründung & 1 & Im gleichschenkligen Dreieck sind beide Basiswinkel gleich; die Innenwinkelsumme ist $180^\circ$: $180^\circ - 2 \cdot 64^\circ$.\\
(a) Ergebnis & 1\E & Spitzenwinkel $52^\circ$\\
(b) $s_a$ & 1 & Die Strecke von der Ecke $A$ zur Mitte der gegenüberliegenden Seite $a$ (gleichwertig: «verbindet $A$ mit dem Mittelpunkt von $\overline{BC}$»).\\
\end{raster}
\fehler{(a) $180^\circ - 64^\circ = 116^\circ$ (nur ein Basiswinkel abgezogen): Ergebniszeile 0, Begründung nur, wenn «beide Basiswinkel gleich» dasteht $\to$ höchstens 2 von 3\,P.
(b) Verwechslung mit der Winkelhalbierenden («halbiert den Winkel bei $A$») oder der Höhe («Lot von $A$ auf $a$»): Zeile (b) 0.
(b) «Mitte der Seite $a$» ohne Ecke $A$ (das wäre die Mittelsenkrechte, wenn senkrecht): Zeile (b) 0.}
\end{lpaufgabe}

\begin{lpaufgabe}{G2}{Dreiecksfläche und Höhe}{4 P}
\begin{raster}
(a) Ansatz & 1 & $A = \tfrac12 \cdot c \cdot h_c$ — Grundseite und \emph{zugehörige} Höhe\\
(a) Ergebnis & 1\E & $A = \tfrac12 \cdot 9 \cdot 4.2 = 18.9\,\text{cm}^2$\\
(b) Ansatz & 1 & Dieselbe Fläche mit der Grundseite $b$: $A = \tfrac12 \cdot b \cdot h_b \Rightarrow h_b = \tfrac{2A}{b}$\\
(b) Ergebnis & 1\E & $h_b = \tfrac{2 \cdot 18.9}{7} = 5.4\,\text{cm}$. Folgefehler aus (a): richtig mit dem eigenen $A$ gerechnet gibt die Ansatzzeile, nicht den (E)-Punkt.\\
\end{raster}
\fehler{Seite statt Höhe, z.\,B. $A = \tfrac12 \cdot 9 \cdot 7 = 31.5\,\text{cm}^2$: (a) beide Zeilen 0.
Den Faktor $\tfrac12$ vergessen ($A = 37.8\,\text{cm}^2$): (a) Ergebniszeile 0; (b) damit $h_b = 10.8$ — Ansatz als Folgefehler 1, Ergebnis 0 $\to$ 2 von 4\,P.
(b) $h_b = \tfrac{A}{b} = 2.7\,\text{cm}$ (Faktor 2 vergessen): (b) beide Zeilen 0.
(b) $h_b = h_c = 4.2\,\text{cm}$ («die Höhe ist dieselbe»): (b) beide Zeilen 0.}
\end{lpaufgabe}

\begin{lpaufgabe}{G3}{Gleichschenkliges Trapez}{5 P}
\begin{raster}
Fläche & 1\E & $A = \tfrac12(a + c) \cdot h = \tfrac12 \cdot 16 \cdot 6 = 48\,\text{cm}^2$\\
Mittellinie & 1\E & $m = \tfrac12(a + c) = 8\,\text{cm}$\\
Überstand & 1 & Gleichschenklig: auf jeder Seite $\tfrac{a - c}{2} = 3\,\text{cm}$; rechtwinkliges Dreieck mit den Katheten $3$ und $6$ (Skizze oder Text).\\
Schenkel & 1\E & $s = \sqrt{3^2 + 6^2} = \sqrt{45} \approx 6.71\,\text{cm}$\\
Umfang & 1\E & $U = 11 + 5 + 2 \cdot 6.71 \approx 29.42\,\text{cm}$ (mit $\sqrt{45}$ exakt: $29.416\ldots$; $29.42$ oder $29.41$ zählt)\\
\end{raster}
\fehler{$A = (a + c) \cdot h = 96\,\text{cm}^2$ (Faktor $\tfrac12$ vergessen): Flächenzeile 0.
Überstand $a - c = 6\,\text{cm}$ statt $3\,\text{cm}$: Überstand- und Schenkelzeile 0 ($s = \sqrt{72} \approx 8.49$); Umfang mit dem eigenen $s$ ($\approx 32.97$) als Folgefehler — (E), also 0 $\to$ 2 von 5\,P.
Die Höhe als Schenkel genommen ($U = 28\,\text{cm}$): Überstand-, Schenkel- und Umfangszeile 0 $\to$ 2 von 5\,P.
Pythagoras falsch angesetzt, $s^2 = 6^2 - 3^2$: Schenkelzeile 0, Umfang 0.}
\end{lpaufgabe}

\begin{lpaufgabe}{G4}{Kreissektor}{4 P}
\begin{raster}
Bogen-Ansatz & 1 & Anteil am Vollkreis: $b = \tfrac{135^\circ}{360^\circ} \cdot 2\pi r$ (gleichwertig: $0.375 \cdot U$ oder $\tfrac38 \cdot 2\pi \cdot 8$)\\
Bogenlänge & 1\E & $b = 6\pi \approx 18.85\,\text{cm}$\\
Sektor-Ansatz & 1 & $A_S = \tfrac{135^\circ}{360^\circ} \cdot \pi r^2$ (gleichwertig: $\tfrac12 \cdot b \cdot r$)\\
Sektorfläche & 1\E & $A_S = 24\pi \approx 75.40\,\text{cm}^2$\\
\end{raster}
\fehler{$r = 8$ als Durchmesser genommen ($r = 4$): $b \approx 9.42\,\text{cm}$, $A_S \approx 18.85\,\text{cm}^2$ — beide Ergebniszeilen 0, Ansätze voll $\to$ 2 von 4\,P.
Sektor mit der Umfangsformel ($\tfrac38 \cdot 2\pi r = 18.85$ als Fläche): Sektor-Ansatz und Ergebnis 0.
Anteil falsch, z.\,B. $\tfrac{135}{180}$: beide Ansatzzeilen 0, Ergebnisse 0.}
\end{lpaufgabe}

\begin{lpaufgabe}{G5}{Kreissegment}{3 P}
\begin{raster}
Sektor & 1\E & $A_S = \tfrac{60^\circ}{360^\circ} \cdot \pi \cdot 10^2 = \tfrac{50\pi}{3} \approx 52.36\,\text{cm}^2$\\
Dreieck & 1 & Bei $\varphi = 60^\circ$ und zwei Radien als Seiten ist das Dreieck gleichseitig (Seite $10\,\text{cm}$): $A_D = \tfrac{\sqrt3}{4} \cdot 10^2 \approx 43.30\,\text{cm}^2$. Gleichwertig: Höhe $\sqrt{100 - 25} \approx 8.66\,\text{cm}$, $A_D = \tfrac12 \cdot 10 \cdot 8.66$.\\
Segment & 1\E & $A = A_S - A_D \approx 9.06\,\text{cm}^2$\\
\end{raster}
\fehler{Dreieck als $\tfrac12 r^2 = 50\,\text{cm}^2$ (wie bei einem rechten Winkel): Dreieckszeile 0; Segment $\approx 2.36$ als Folgefehler — (E), also 0 $\to$ 1 von 3\,P.
Sektor plus statt minus Dreieck: Segmentzeile 0.
Ganzer Kreis minus Dreieck: Sektor- und Segmentzeile 0.}
\end{lpaufgabe}

\begin{lpaufgabe}{G6}{Schatten}{3 P}
\begin{raster}
Skizze und Ansatz & 1 & Gleicher Sonnenstand: zwei ähnliche rechtwinklige Dreiecke (Höhe und Schatten als Katheten). $\tfrac{h}{14} = \tfrac{1.8}{2.4}$ oder gleichwertig $k = \tfrac{14}{2.4}$.\\
Verhältnis & 1 & $\tfrac{1.8}{2.4} = 0.75$ bzw. $k \approx 5.83$, entsprechende Seiten richtig zugeordnet\\
Ergebnis & 1\E & $h = 0.75 \cdot 14 = 10.5\,\text{m}$\\
\end{raster}
\fehler{Verhältnis falsch herum, $h = \tfrac{2.4}{1.8} \cdot 14 \approx 18.67\,\text{m}$: Verhältnis- und Ergebniszeile 0 $\to$ 1 von 3\,P (Plausibilität: der Mast muss kürzer sein als sein Schatten, wie die Person).
Differenzen statt Verhältnis ($14 - 2.4 + 1.8$): alle Zeilen 0.
Ohne Skizze, aber mit richtigem Ansatz: Ansatzzeile trotzdem 1 — die Skizze wird verlangt, ihr Fehlen kostet aber keinen eigenen Punkt.}
\end{lpaufgabe}

\begin{lpaufgabe}{G7}{Massstab und Fläche}{3 P}
\begin{raster}
Längenfaktor & 1 & Massstab $1 : 50$: alle Längen in Wirklichkeit $50$-mal so gross, $k = 50$\\
Flächenfaktor & 1 & Flächen mit $k^2 = 2500$: $120\,\text{cm}^2 \cdot 2500 = 300\,000\,\text{cm}^2$\\
Ergebnis in m² & 1\E & $1\,\text{m}^2 = 10\,000\,\text{cm}^2$: $300\,000\,\text{cm}^2 = 30\,\text{m}^2$\\
\end{raster}
\fehler{Flächenfaktor $k$ statt $k^2$ (Massstab nur linear): $120 \cdot 50 = 6000\,\text{cm}^2 = 0.6\,\text{m}^2$ — Flächenfaktor- und Ergebniszeile 0 $\to$ 1 von 3\,P.
Umrechnung mit $1\,\text{m}^2 = 100\,\text{cm}^2$ ($3000\,\text{m}^2$): Ergebniszeile 0.
Ergebnis nur in $\text{cm}^2$: Ergebniszeile 0 — die Aufgabe verlangt $\text{m}^2$.}
\end{lpaufgabe}

\needspace{14\baselineskip}
\section*{4 · Selbsteinschätzung}
\begin{tabularx}{\linewidth}{@{}l X@{}}\toprule
22 – 25 P & Die geprüften Teile sitzen. Wo du Punkte verloren hast: das Kapitel dieser Aufgabe nochmals (Zuordnung unten).\\
17 – 21 P & Den schwächsten Teil nochmals: Tüfteln und Übungen des Kapitels, in dem du die meisten Punkte verloren hast.\\
11 – 16 P & Zurück zu den Kapiteln aller Aufgaben, in denen du Punkte verloren hast.\\
0 – 10 P & Zurück zu Kapitel 1 und von dort der Reihe nach weiter.\\\midrule
\multicolumn{2}{@{}p{\linewidth}@{}}{\small\textbf{Aufgabe $\to$ Kapitel:} G1 $\to$ Kapitel 1 (Dreiecke beschreiben); G2 $\to$ Kapitel 2 (Dreiecksfläche und Höhe); G3 $\to$ Kapitel 3 (Vierecke); G4, G5 $\to$ Kapitel 4 (Kreis und Kreisteile); G6, G7 $\to$ Kapitel 5 (Ähnlichkeit)}\\\bottomrule
\end{tabularx}
\end{document}
